\chapter{Inférence à faible latence}\label{ch:ml:low-latency-inference}

Un modèle d'arbres boostés par gradient à 300 arbres répond en 265~µs via l'interface Python de sa bibliothèque et en 2.6~µs
sous forme de C++ généré, mesurés sur le même ordinateur portable~; le modèle n'a pas changé, et les deux réponses sont égales au
dernier bit. La plus grande part du premier chiffre ne tient pas au modèle~: elle vient de l'interpréteur, de l'enveloppe, de la conversion
d'une ligne en tableau et du retour. Ce chapitre porte sur la sortie des modèles hors de cet environnement~: compiler les arbres en code, réduire les petits
réseaux à des entiers sur huit bits, décider quand regrouper les requêtes, et le faire en C++20 et en Rust avec les résultats exacts de la référence
Python. Il s'appuie sur la méthode de mesure de la latence du livre~13, et ses temps sont ceux d'un seul ordinateur portable, signalés comme tels~;
ce que les tests vérifient est la parité, et ce que le chapitre conclut porte sur des ordres de grandeur.

\section{Le budget d'inférence}

\begin{definition}[Latence d'inférence]\label{def:ml:low-latency-inference:latency}
La \emph{latence d'inférence}\index{latence d'inférence} désigne le temps écoulé entre le moment où les entrées d'un modèle sont disponibles et le moment
où sa sortie l'est, mesuré à chaque appel et rapporté sous forme de distribution (médiane et queue), dans le budget de latence de la stratégie
(livre~13, chapitre~1).
\end{definition}

Un modèle de market making ou d'exécution décide dans un budget de quelques microsecondes à quelques centaines, partagé avec
le calcul des variables (chapitre~24), les contrôles de risque et le réseau. L'inférence doit tenir sa part dans la queue de la distribution, et non dans la médiane~: la décision
qui arrive en retard est celle qui est prise lorsque le marché a bougé. Les deux modèles du chapitre sont ceux
de la fixture de \texttt{firm.mlinfer}, ajustés à un problème synthétique à 16 variables~: une forêt LightGBM (chapitre~5) de 300 arbres à 15 feuilles chacun, soit 4\,200
coupures au total, et un \omterm{def:ml:neural-networks-for-noisy-tabular-data:mlp}{perceptron multicouche} (chapitre~7) de $16\to32\to16\to1$ unités. Sur 20\,000 nouvelles observations, le coefficient d'information de rang de la forêt
(rank IC) vaut 0.224 et celui du réseau 0.182.

\section{Compiler les arbres}

\begin{definition}[Compilation de modèle]\label{def:ml:low-latency-inference:compile}
La \emph{compilation de modèle}\index{compilation de modèle} traduit un modèle entraîné en code ou en structures de données spécialisées pour
l'inférence sur une cible donnée (code source généré, disposition du tableau à plat, description matérielle), en retirant le framework d'entraînement du
chemin de prédiction (Asadi, Lin et de Vries, 2014~; Lucchese et coauteurs, 2015, pour les forêts de classement).
\end{definition}

La forêt est exportée de deux façons. Sous forme de tableaux plats d'indices de variables, de seuils et d'indices d'enfants, une entrée par nœud, que parcourt
une boucle (\cref{lst:ml:inf:loop})~: la même boucle tourne en Python, en C++20 et en Rust. Et sous forme de C++ généré, avec un
\texttt{if} imbriqué par coupure et les seuils écrits en littéraux (\cref{lst:ml:inf:gen}), soit 17\,107 lignes que le compilateur transforme en
branchements droits. Les deux reproduisent exactement les prédictions de LightGBM sur les 200 vecteurs de test communs, car elles visitent les mêmes
feuilles et ajoutent les valeurs des feuilles dans le même ordre en double précision~; les seuils survivent au passage par le texte parce qu'ils sont
écrits dans la représentation la plus courte qui se relit comme le même double.

La \Cref{fig:ml:inf:latency} montre les latences mesurées. L'interface scikit-learn de LightGBM prend 265~µs par appel à une seule
ligne en médiane et 544~µs au 99\up{e}~centile~; la même boucle plate en Python n'est pas plus rapide (300~µs), car le coût de l'interpréteur
par nœud domine. En C++, la boucle prend 6.2~µs et les branchements générés 2.6~µs (5.7~µs au 99\up{e}~centile), soit une
centaine de fois moins que l'appel de la bibliothèque, aussi bien en médiane que dans la queue~; la boucle en Rust prend 7.8~µs. Les branchements l'emportent sur
la boucle, parce que les seuils et les adresses des enfants sont dans le flux d'instructions plutôt qu'en mémoire, et que le compilateur
peut disposer le chemin le plus probable~; avec 300 arbres et 4\,200 coupures, la forêt entière tient encore dans les caches (livre~13, chapitre~3).

\begin{omfigure}
\begin{tikzpicture}
\begin{semilogxaxis}[width=11cm,height=7.5cm,xbar,bar width=5pt,bar shift=0pt,xmin=10,xmax=2e6,
  xlabel={latence par appel (nanosecondes)},ytick={0,1,2,3,4,5,6,7,8,9},
  yticklabels={LightGBM predict (Python),forêt exportée en Python,MLP float en PyTorch,MLP int8 en NumPy,forêt C++ (boucle),forêt C++ (branches),MLP int8 C++,surcoût de l'horloge,forêt Rust,MLP int8 Rust},y dir=reverse,
  tick label style={font=\small},label style={font=\small},grid=major,grid style={gray!25},area legend,
  legend style={font=\footnotesize,at={(0.5,-0.2)},anchor=north,draw=none,fill=none,/tikz/every even column/.append style={column sep=8pt}},legend columns=2,
  table/col sep=comma,enlarge y limits=0.06,log origin=infty]
\addplot[fill=omThm,draw=none,bar shift=2.6pt] table[x=p99_ns,y=i]{figdata/ml/26-low-latency-inference/measured_latency.csv};
\addplot[fill=omDef,draw=none,bar shift=-2.6pt] table[x=median_ns,y=i]{figdata/ml/26-low-latency-inference/measured_latency.csv};
\legend{99\up{e}~centile,médiane}
\end{semilogxaxis}
\end{tikzpicture}

\omcaption{Latence d'un appel isolé pour les deux mêmes modèles, par voie d'exécution~; 200\,000 appels en C++ et en Rust, 5\,000 en Python (300
pour la forêt en Python pur), chronométrés un par un. Mesurés sur un ordinateur portable (Intel Core Ultra 7 155H) sous WSL2 en basse priorité, sans
cœurs isolés~; d'autres tâches ont pu tourner. Données~: \texttt{bench\_infer.py}, \texttt{measured\_latency.csv}.}\label{fig:ml:inf:latency}
\end{omfigure}

\section{Petits réseaux~: disposition, fusion et quantification}

\begin{definition}[Fusion d'opérateurs, quantification, quantification post-entraînement, entraînement conscient de la quantification]\label{def:ml:low-latency-inference:quant}
La \emph{fusion d'opérateurs}\index{fusion d'opérateurs} calcule plusieurs opérations consécutives (un produit matriciel, un biais, une
activation, une remise à l'échelle) en un seul passage sur les données. La \emph{quantification}\index{quantification} représente les poids et
les activations par de petits entiers munis de facteurs d'échelle, de sorte que l'inférence s'exécute en arithmétique entière (Jacob et coauteurs, 2018).
La \emph{quantification post-entraînement}\index{quantification post-entraînement} calcule les entiers et les échelles à partir d'un modèle flottant entraîné
et d'un échantillon de calibration~; l'\emph{entraînement conscient de la quantification}\index{entraînement conscient de la quantification} fait un \omterm{def:ml:representation-learning:ssl}{ajustement fin} du modèle en simulant
l'arrondi dans la passe avant, afin qu'il apprenne des poids qui y résistent (Nagel et coauteurs, 2021).
\end{definition}

Le réseau du chapitre est quantifié de façon symétrique, avec une échelle par tenseur~: poids et activations en entiers de $[-127, 127]$,
biais en entiers sur 32~bits à la valeur du produit des échelles d'entrée et de poids, et entre les couches un multiplicateur à virgule fixe et un décalage
qui ramènent l'accumulateur sur 32~bits à huit bits, en arrondissant à la moitié en s'éloignant de zéro (\cref{lst:ml:inf:mlp}). Chaque couche est un
seul passage fusionné~: produits, biais, remise à l'échelle et ReLU. L'arithmétique entière rend la parité exacte par construction~: la référence NumPy
le noyau C++ et le noyau Rust s'accordent sur chacun des 200 vecteurs de test, au bit près. Les poids sur huit bits
occupent 1\,236 octets contre 4\,356 en float32, et un appel prend 0.47~µs en C++ et 0.58~µs en Rust.

\begin{table}[htbp]
\centering\small
\begin{tabular}{@{}lccc@{}}
\toprule
& \multicolumn{3}{c}{bits par poids et par activation}\\
\cmidrule(l){2-4}
coefficient d'information de rang sur 20\,000 nouvelles observations & 8 & 6 & 4\\
\midrule
quantifié après entraînement & 0.182 & 0.181 & 0.143\\
\omterm{def:ml:representation-learning:ssl}{ajustement fin} conscient de la \omterm{def:ml:low-latency-inference:quant}{quantification}, puis quantifié & 0.179 & 0.176 & 0.164\\
\bottomrule
\end{tabular}
\caption{Le coefficient d'information du réseau après \omterm{def:ml:low-latency-inference:quant}{quantification} (flottant~: 0.182~; la forêt~: 0.224). Données~:
\texttt{ml\_infer.accuracy}.}\label{tab:ml:inf:quant}
\end{table}

Le \Cref{tab:ml:inf:quant} donne la seconde moitié du résultat chiffré du chapitre. À huit bits, la \omterm{def:ml:low-latency-inference:quant}{quantification} après entraînement ne coûte
rien de mesurable (0.1817 contre 0.1822~; les sorties du modèle entier diffèrent de celles du modèle flottant de 4.5~\% de leur
écart type, et le classement bouge à peine)~; l'\omterm{def:ml:representation-learning:ssl}{ajustement fin} conscient de la \omterm{def:ml:low-latency-inference:quant}{quantification}, qui consacre dix \omterm{def:ml:neural-networks-for-noisy-tabular-data:backprop}{époques} à s'adapter à l'arrondi,
perd un peu (0.179), car l'\omterm{def:ml:representation-learning:ssl}{ajustement fin} lui-même déplace le modèle. À six bits, le constat est le même. À quatre bits, la \omterm{def:ml:low-latency-inference:quant}{quantification post-entraînement}
perd un cinquième du signal (0.143), et l'\omterm{def:ml:low-latency-inference:quant}{entraînement conscient de la quantification} en récupère plus de la moitié de la perte (0.164).
Huit bits sont gratuits pour un réseau aussi petit~; c'est en dessous que l'entraînement conscient de l'arrondi se rentabilise.

\section{Regrouper contre latence}

\begin{definition}[Regroupement dynamique]\label{def:ml:low-latency-inference:batch}
Le \emph{regroupement dynamique}\index{regroupement dynamique} (dynamic batching) rassemble les requêtes qui arrivent à peu de distance les unes des autres en un seul appel du modèle, en échangeant
le temps d'attente de la première requête contre un coût moindre par requête.
\end{definition}

Pour une bibliothèque dont le surcoût par appel domine le coût par ligne, le regroupement est un levier important (la \cref{fig:ml:inf:batch})~: un appel de LightGBM
prend 273~µs pour une seule ligne et 3.1~ms pour 1\,024~: le coût par ligne tombe ainsi de 273~µs à 3.0~µs. Un travail de recherche qui évalue
un univers à chaque seconde doit regrouper ses appels. Une décision de trading sur un seul instrument ne peut pas attendre l'arrivée des requêtes des autres, et un
modèle compilé coûte déjà quelques microsecondes par ligne~: là, le regroupement achète un débit que la stratégie n'utilise pas, au
prix d'une latence qu'elle ne peut pas se permettre. Quand on regroupe, la latence se mesure à partir de l'arrivée de chaque requête, et non
à partir du début du lot, sans quoi le temps d'attente disparaît des statistiques (omission coordonnée, livre~13, chapitre~5).

\begin{omfigure}
\begin{tikzpicture}
\begin{loglogaxis}[width=10.5cm,height=6cm,xlabel={lignes par appel},ylabel={microsecondes},tick label style={font=\small},
  label style={font=\small},xmin=0.8,xmax=1300,grid=major,grid style={gray!25},legend cell align=left,
  legend style={font=\footnotesize,at={(0.02,0.02)},anchor=south west,fill=white,draw=none},table/col sep=comma]
\addplot[omThm,very thick,mark=*] table[x=batch,y=call_us]{figdata/ml/26-low-latency-inference/measured_batch.csv};
\addplot[omDef,very thick,mark=square*] table[x=batch,y=per_row_us]{figdata/ml/26-low-latency-inference/measured_batch.csv};
\legend{par appel,par ligne}
\end{loglogaxis}
\end{tikzpicture}

\omcaption{Latence de prédiction de LightGBM (médiane) en fonction de la taille du lot, via son interface Python, par appel et par ligne.
Même machine et mêmes conditions que la \cref{fig:ml:inf:latency}. Données~: \texttt{bench\_infer.py},
\texttt{measured\_batch.csv}.}\label{fig:ml:inf:batch}
\end{omfigure}

\section{Inférence en C++, en Rust et sur matériel programmable}

Les noyaux chargent leur modèle une seule fois, en allouant de la mémoire, puis prédisent sans toucher au tas~: des tampons de taille fixe sur la pile, sans
exceptions, sans appels virtuels, les tableaux du modèle en lecture seule et partagés. La version Rust (le \cref{lst:ml:inf:rust}) suit le même
algorithme avec des vérifications de bornes, qui lui coûtent environ un quart par rapport au C++ ici~; les deux pourraient les supprimer par une indexation non vérifiée
une fois les tableaux du modèle validés au chargement. Au-delà du CPU, le matériel programmable exécute arbres et petits réseaux en
un nombre fixe de cycles d'horloge, disposés sous forme de circuit~: Duarte et coauteurs (2018) ont compilé de petits réseaux pour des FPGA dans les
systèmes de déclenchement de la physique des particules, où les décisions disposent de budgets de latence fixes du même ordre. La même discipline de
parité s'applique~: un modèle matériel est validé contre la référence Python sur des vecteurs communs, bit à bit lorsque l'arithmétique est
entière.

\begin{method}[Mettre un modèle sur le chemin rapide]\label{met:ml:inf:method}
\begin{enumerate}
\item Mesurer la médiane et la queue de l'appel de la bibliothèque sur la machine cible, une requête à la fois, avant de rien modifier.
\item Exporter le modèle sous une forme portable et le compiler (arbres en code ou en tableaux plats~; réseaux en noyaux entiers fusionnés),
sans allocation sur le hot path.
\item Tester la parité sur des vecteurs communs~: exacte pour les arbres et les réseaux entiers, dans les tolérances énoncées pour les
réseaux en virgule flottante.
\item Quantifier après l'entraînement et mesurer la perte de signal~; ne recourir à l'\omterm{def:ml:low-latency-inference:quant}{entraînement conscient de la quantification} que lorsque la perte importe.
\item Enregistrer l'artefact compilé avec son rapport de parité à côté du modèle dont il provient (chapitre~25).
\end{enumerate}
\end{method}

\section{Tutoriel~: deux microsecondes pour une forêt}\label{tut:ml:low-latency-inference:tut}

\begin{tutorial}
\textbf{Objectif.} Exporter une forêt et un réseau, les compiler en C++20 et en Rust, vérifier la parité, quantifier et mesurer.
\textbf{État final~:} les \cref{fig:ml:inf:latency,fig:ml:inf:batch}, le \cref{tab:ml:inf:quant}.
\begin{tutsteps}
\item \textbf{Arbres sous forme de code généré.}
\omcode{code/firm/mlinfer/firm_mlinfer.py}{84}{103}{Génération des branches C++ imbriquées à partir de la forêt exportée.}\label{lst:ml:inf:gen}
\item \textbf{La forêt à plat en C++20.}
\omcode{code/firm/mlinfer/cpp/mlinfer.hpp}{20}{27}{Parcours des tableaux plats de nœuds.}\label{lst:ml:inf:loop}
\item \textbf{Le réseau int8 en C++20.}
\omcode{code/firm/mlinfer/cpp/mlinfer.hpp}{66}{94}{Inférence entière avec requantification en virgule fixe.}\label{lst:ml:inf:mlp}
\item \textbf{La forêt à plat en Rust.}
\omcode{code/firm/mlinfer/rust/src/lib.rs}{15}{26}{La même boucle en Rust.}\label{lst:ml:inf:rust}
\item \textbf{Exécuter} \texttt{make\_mlinfer\_fixture.py} (la fixture), les tests de parité C++ et Rust (\texttt{make test-code}),
\texttt{ml\_infer.accuracy()}, et, une seule fois, \texttt{bench\_infer.py} sous \texttt{nice}.
\end{tutsteps}
\textbf{Pour aller plus loin.} Stocker les nœuds de la forêt en ordre de parcours en largeur pour améliorer la localité, ou évaluer plusieurs arbres
à la fois avec des instructions vectorielles (livre~13, chapitre~14), puis mesurer.
\end{tutorial}

\section{Construction~: les noyaux d'inférence}\label{bld:ml:low-latency-inference:mlinfer}

\begin{build}
\bfield{Rôle} Des modèles sur le chemin rapide, avec les résultats exacts de la référence Python.
\bfield{Interface} Python~: \texttt{export\_forest}, \texttt{forest\_predict}, \texttt{write\_forest}, \texttt{read\_forest},
\texttt{cpp\_branches}, \texttt{MLP}, \texttt{fit\_mlp(model, X, y, epochs, seed, qat, lr, batch, bits)},
\texttt{quantise(model, X\_calib, bits)}, \texttt{int8\_forward}, \texttt{write\_mlp}~; \texttt{make\_mlinfer\_fixture.py}. C++20~: \texttt{mlinfer::Forest},
\texttt{read\_forest}, \texttt{Mlp}, \texttt{read\_mlp}, les \texttt{forest\_branches} générées. Rust~: \texttt{Forest},
\texttt{Mlp}, \texttt{requant}.
\bfield{Règles} Aucune allocation, aucune exception ni aucun appel virtuel sur le chemin de prédiction~; la parité sur des vecteurs communs est exacte~; les durées
sont mesurées, étiquetées avec la machine, et ne sont jamais affirmées.
\bfield{Tests d'acceptation} \texttt{code/firm/mlinfer/tests/} (Python~: l'export égale LightGBM, la fixture est reproductible, la référence entière
contre le modèle en virgule flottante), \texttt{cpp/mlinfer\_test.cpp} et \texttt{rust} (\texttt{cargo test})~: parité exacte de la
forêt à plat, des branches générées et du réseau int8 sur les 200 vecteurs.
\bfield{Pour aller plus loin} Évaluation SIMD de plusieurs arbres~; \omterm{def:ml:low-latency-inference:quant}{quantification} par canal~; un export en description matérielle du réseau int8.
\end{build}

\begin{omsources}
\item B.~Jacob et coauteurs, «~Quantization and training of neural networks for efficient integer-arithmetic-only
inference~», arXiv:1712.05877, 2018.
\item M.~Nagel et coauteurs, «~A white paper on neural network quantization~», arXiv:2106.08295, 2021.
\item N.~Asadi, J.~Lin et A.~P.~de Vries, «~Runtime optimizations for tree-based machine learning models~», \emph{IEEE
Transactions on Knowledge and Data Engineering}, 2014.
\item C.~Lucchese et coauteurs, «~QuickScorer: a fast algorithm to rank documents with additive ensembles of regression
trees~», \emph{SIGIR}, 2015.
\item J.~Duarte et coauteurs, «~Fast inference of deep neural networks in FPGAs for particle physics~», arXiv:1804.06913,
2018.
\end{omsources}

\section{Exercices}

\begin{exercise}[$\star$]\label{exo:ml:low-latency-inference:1}
Un tenseur int8 symétrique a pour échelle $s = \max|w|/127$. Calculer la plus grande erreur d'arrondi d'un poids, en proportion de $\max|w|$.
\end{exercise}

\begin{exercise}[$\star$]\label{exo:ml:low-latency-inference:2}
D'après la \cref{fig:ml:inf:latency}, calculer le rapport entre la latence au 99\up{e}~centile de la forêt C++ générée et celle de l'appel de la bibliothèque.
\end{exercise}

\begin{exercise}[$\star$]\label{exo:ml:low-latency-inference:3}
Expliquer pourquoi la forêt compilée peut reproduire exactement la prédiction de LightGBM, alors qu'un réseau flottant compilé ne peut
généralement pas reproduire celle de PyTorch.
\end{exercise}

\begin{exercise}[$\star\star$]\label{exo:ml:low-latency-inference:4}
Écrire le multiplicateur et le décalage en virgule fixe pour un multiplicateur réel de 0.0123. Pourquoi utiliser des entiers pour cela~?
\end{exercise}

\begin{exercise}[$\star\star$]\label{exo:ml:low-latency-inference:5}
Pourquoi l'\omterm{def:ml:low-latency-inference:quant}{entraînement conscient de la quantification} a-t-il perdu un peu à huit bits et gagné à quatre~?
\end{exercise}

\begin{exercise}[$\star\star$]\label{exo:ml:low-latency-inference:6}
\emph{Trouver la faille.} «~La latence médiane de notre modèle compilé est de 800 nanosecondes, mesurée en chronométrant 10\,000 prédictions en un seul
lot et en divisant par 10\,000.~»
\end{exercise}

\begin{exercise}[$\star\star\star$]\label{exo:ml:low-latency-inference:7}
\emph{Programmation.} Quantifier le réseau sur quatre bits après l'entraînement, avec une échelle de poids par unité de sortie au lieu d'une par
tenseur. Quelle part de la perte par rapport au modèle flottant cela permet-il de récupérer~?
\end{exercise}

\begin{exercise}[$\star\star\star$]\label{exo:ml:low-latency-inference:8}
Avec des échelles symétriques $s_x$, $s_w$ et $s_y$ pour l'entrée, les poids et la sortie d'une couche, écrire la sortie flottante de la couche
en fonction des entiers $X_i$, $W_i$ et d'un biais entier $B$, puis en déduire l'échelle à laquelle $B$ doit être stocké ainsi que le multiplicateur qui fait passer
l'accumulateur aux entiers de sortie.
\end{exercise}

\section{Problème~: Deux microsecondes pour une forêt}

\begin{problem}[{Problème du week-end --- le modèle n'est pas la latence}]\label{pb:ml:low-latency-inference:1}
La forêt, le réseau, les noyaux et les mesures du chapitre.

\medskip\noindent\textbf{Partie I --- Les modèles.}
\begin{enumerate}
\item Quels sont les deux modèles, et quelle est leur qualité~?
\item Comment la forêt est-elle exportée, et sous combien de formes~?
\item Pourquoi les noyaux sont-ils exacts~?
\item Quelle est la taille en mémoire de chaque modèle~?
\end{enumerate}

\medskip\noindent\textbf{Partie II --- Latence.}
\begin{enumerate}[resume]
\item Donner les latences médiane et au 99\up{e}~centile de chaque chemin.
\item Où part le temps de l'appel de la bibliothèque~?
\item Pourquoi les branches générées battent-elles la boucle~?
\item Que coûte la version Rust face au C++, et pourquoi~?
\end{enumerate}

\medskip\noindent\textbf{Partie III --- \omterm{def:ml:low-latency-inference:quant}{Quantification} et regroupement.}
\begin{enumerate}[resume]
\item Que coûte l'int8 en coefficient d'information, après entraînement et avec l'\omterm{def:ml:low-latency-inference:quant}{entraînement conscient de la quantification}~?
\item Que change-t-il à quatre bits~?
\item Comment la latence de LightGBM évolue-t-elle avec la taille du lot~?
\item Quand faudrait-il regrouper les requêtes~?
\end{enumerate}

\medskip\noindent\textbf{Partie IV --- Le verdict.}
\begin{enumerate}[resume]
\item Énoncer le \emph{résultat nommé}~: le rapport entre la latence de queue du modèle compilé et celle de l'appel de la bibliothèque, ainsi que le coefficient d'information perdu
à cause de la \omterm{def:ml:low-latency-inference:quant}{quantification} int8, avec et sans \omterm{def:ml:low-latency-inference:quant}{entraînement conscient de la quantification}.
\item Quel modèle déploieriez-vous dans un budget de cinq microsecondes, et sous quelle forme~?
\item Que doit contenir le registre pour un modèle compilé~?
\item Comment rendre ces mesures fiables sur une machine de production~?
\item Que changerait-il avec 5\,000 arbres~?
\item Où un matériel programmable vaudrait-il son coût~?
\item Comment un biais de \omterm{def:ml:data-and-feature-stores:store}{feature store} (chapitre~24) se manifesterait-il à cette couche~?
\item En une phrase~: où part le temps d'inférence~?
\end{enumerate}
\end{problem}

\section{Questions d'entretien}

\begin{interviewq}[$\star$ \iqroles{développeur, ingénieur ML}]\label{iq:ml:low-latency-inference:1}
Comment feriez-vous pour qu'un modèle à renforcement de gradient prédise en quelques microsecondes~?
\end{interviewq}

\begin{interviewq}[$\star\star$ \iqroles{ingénieur ML}]\label{iq:ml:low-latency-inference:2}
Expliquez la \omterm{def:ml:low-latency-inference:quant}{quantification} int8~: échelles, points zéro, accumulateurs et requantification.
\end{interviewq}

\begin{interviewq}[$\star\star$ \iqroles{développeur}]\label{iq:ml:low-latency-inference:3}
Quelles règles imposeriez-vous au code du chemin d'inférence critique~?
\end{interviewq}

\begin{interviewq}[$\star\star$ \iqroles{ingénieur ML, développeur}]\label{iq:ml:low-latency-inference:4}
Comment vérifiez-vous qu'un modèle compilé est bien celui que vous avez entraîné~?
\end{interviewq}

\begin{interviewq}[$\star\star$ \iqroles{développeur}]\label{iq:ml:low-latency-inference:5}
Quand le regroupement améliore-t-il la \omterm{def:ml:low-latency-inference:latency}{latence d'inférence}, et quand la dégrade-t-il~?
\end{interviewq}

\begin{interviewq}[$\star\star\star$ \iqroles{ingénieur ML}]\label{iq:ml:low-latency-inference:6}
\omterm{def:ml:low-latency-inference:quant}{Quantification post-entraînement} ou \omterm{def:ml:low-latency-inference:quant}{entraînement conscient de la quantification}~? Comment trancheriez-vous~?
\end{interviewq}
